\documentclass[12pt]{article}
\usepackage{amsmath}
\usepackage{physics}

\begin{document}

\section*{Thermal Efficiency Derivation for the Cycle}

The thermal efficiency $\eta$ of a thermodynamic cycle is defined as the ratio of net work to heat input:
\begin{equation}
\eta = 1 - \frac{Q_{\text{out}}}{Q_{\text{in}}}.
\end{equation}

Assume the cycle follows these states in order:

\begin{enumerate}
    \item State 1 $\rightarrow$ 2: Adiabatic expansion ($Q = 0$)
    \item State 2 $\rightarrow$ 3: Isochoric heat removal ($Q_{\text{out,1}}$)
    \item State 3 $\rightarrow$ 4: Isobaric heat removal ($Q_{\text{out,2}}$)
    \item State 4 $\rightarrow$ 5: Isobaric heat addition ($Q_{\text{in,1}}$)
    \item State 5 $\rightarrow$ 1: Isochoric heat addition ($Q_{\text{in,2}}$)
\end{enumerate}

\subsection*{1. Heat Input ($Q_{\text{in}}$)}
Heat is added during the isobaric (4-5) and isochoric (5-1) processes:
\begin{equation}
Q_{\text{in}} = m C_p (T_5 - T_4) + m C_v (T_1 - T_5)
\end{equation}

\subsection*{2. Heat Output ($Q_{\text{out}}$)}
Heat is removed during the isochoric (2-3) and isobaric (3-4) processes:
\begin{equation}
Q_{\text{out}} = m C_v (T_2 - T_3) + m C_p (T_3 - T_4)
\end{equation}

\subsection*{3. Efficiency Formula}
Substituting $C_p = \gamma C_v$ and simplifying:
\begin{align}
\eta = 1 - \frac{m C_v (T_2 - T_3) + m C_p (T_3 - T_4)}{m C_p (T_5 - T_4) + m C_v (T_1 - T_5)} \\
= 1 - \frac{C_v (T_2 - T_3) + \gamma C_v (T_3 - T_4)}{\gamma C_v (T_5 - T_4) + C_v (T_1 - T_5)} \\
= 1 - \frac{(T_2 - T_3) + \gamma (T_3 - T_4)}{(T_1 - T_5) + \gamma (T_5 - T_4)}
\end{align}

\subsection*{Answer}
The final expression for the thermal efficiency of the cycle is:
\begin{equation}
\eta = 1 - \frac{T_2 - T_3 + \gamma (T_3 - T_4)}{T_1 - T_5 + \gamma (T_5 - T_4)}
\end{equation}

\end{document}

